% ============================================================
%  Viyog #215 — Response to Reviewers
%  IEEE TCAD double-column; 4 pages.
%  All manuscript changes referenced here are highlighted in BLUE in the
%  revised PDF (all blue text added/changed for this revision round).
% ============================================================
\documentclass[10pt,letterpaper,journal]{IEEEtran}
\usepackage{lineno}

\usepackage{cite}
\usepackage{amsmath,amssymb}
\usepackage{graphicx}
\usepackage[table]{xcolor}
\usepackage{booktabs}
\usepackage{enumitem}
\usepackage{array}
\usepackage{multirow}
\usepackage{balance}
\usepackage{hyperref}

% ---- Macros -----------------------------------------------------------------
% Reviewer comment block: compact dark-grey italic (tight, no box chrome)
\newcommand{\rc}[1]{\smallskip\noindent\textit{\textcolor{black!75}{#1}}\par\smallskip}

\newcommand{\rr}{\noindent\textbf{Response.}~}
\newcommand{\chg}[1]{\textcolor{blue}{\textbf{[Changed:} #1\textbf{]}}}
\newcommand{\Linf}{$L_\infty$}
\newcommand{\xmark}{\textcolor{red}{\ensuremath{\times}}}
\newcommand{\done}[1]{\textcolor{teal}{\textbf{[Done: #1]}}}

\title{Response to Reviewers --- Paper \#215\\
``Viyog: Separating Adversarial and Out-of-Distribution Inputs\\
via First-Layer Dormant-Band Statistics''}
\author{Anonymous Author(s)}

\begin{document}
\maketitle

\noindent We thank the four reviewers for their thorough, constructive reviews.
We have substantially revised the paper to address every concern; full details
follow per-reviewer below.
\textbf{All revisions are highlighted in blue in the revised PDF.}
Each \textcolor{teal}{\textbf{[Done:]}} tag names the section, figure, or table
where the change appears; PDF page numbers are given in parentheses after each tag.

% ── Quick-reference: what changed ─────────────────────────────────────────────
\section*{Summary of Changes in the Revised Manuscript}

\noindent\textbf{New / updated sections:}
\begin{itemize}[leftmargin=1.4em,nosep,topsep=2pt]
  \item \textbf{Abstract \& Sec.~I (Introduction):} deployment framing front-loaded; $L_\infty$-form $\to$ TV statistic (Viyog) naming clarified; structured TV-mitigation explanation added.
  \item \textbf{Sec.~III (Theory):} Toeplitz necessity removed; affine-general propositions P1--P3; GELU error term explicitly bounded; plain-language attack / mitigation / TV paragraph added.
  \item \textbf{Sec.~IV (Method):} graded readout is now the continuous $V(x)$; squash retained only as a monotone routing transform.
  \item \textbf{Sec.~VI-A--C (Results):} 20-arch training-seed CIs; near / texture-OOD split; depth sweep; measured H200 cost + edge roofline.
  \item \textbf{Sec.~VII (Adaptive attacks):} full A0--A3 ladder + EOT (20 archs + GTSRB).
  \item \textbf{Sec.~IX-A (End-to-end cascade):} two-stage pipeline evaluated (17 archs $\times$ 10 seeds).
  \item \textbf{Sec.~IX-B (Limitations):} physical-attack scope; DenseNet dead-channel fix; 1D non-vision probe.
\end{itemize}

\smallskip
\noindent\textbf{New / updated tables:}
Tab.~I (related work), Tab.~II (coverage), Tab.~III (main results, directionless 20-arch),
Tab.~IV (detector cost), Tab.~V (end-to-end cascade), Tab.~VI (limitations).

\smallskip
\noindent\textbf{New / updated figures:}
Figs.~1--3, 5, 7--15 (main paper); App.~Figs.~16--18.

\smallskip
\noindent\textit{All new and changed text is highlighted in \textcolor{blue}{blue} in the revised PDF.}

% ── Main results summary table ────────────────────────────────────────────────
\begin{table}[t]\centering\footnotesize\setlength{\tabcolsep}{1.5pt}
\caption{Directionless AUROC / FPR@95\% on 20 CIFAR-100 architectures ($\times$ =
direction-inconsistent, not deployable). Viyog is the only method with
deployable T2 and T3, at $0.3$\,KB state.}
\begin{tabular}{@{}lcccc@{}}
\toprule
Method & T1 ID/OOD & T2 ID/ADV & T3 OOD/ADV & State \\\midrule
Energy / MaxLogit / GEN & \textbf{0.85--0.85} & 0.36--0.41\,\xmark & 0.79--0.81 & 4\,B \\
Mahalanobis / KNN & 0.88--0.90 & 0.84--0.88 & 0.69--0.71 & 4.5--40\,MB \\
\rowcolor{teal!12}\textbf{Viyog} (dorm.\ TV) & 0.69 & \textbf{0.966} & \textbf{0.824} & \textbf{0.3\,KB} \\
FirstConv-\Linf\ (old $L_\infty$-form) & 0.60\,\xmark & 0.56\,\xmark & 0.45\,\xmark & 8\,B \\
\bottomrule
\end{tabular}
\end{table}


% ═══════════════════════════════════════════════════════════════════════════════
\section*{Reviewer A}
\noindent\textit{We address all 7 detailed comments
(DC~1--7) below. Weaknesses W1--W5 map to DC~1--5 respectively.}

% ------ DC 1 ------------------------------------------------------------------
\rc{DC~1. ``The paper repeatedly emphasizes the fact that Viyog is a second stage
detector that operates on the inputs flagged as non-ID by some first-stage
detectors. However, has an end-to-end evaluation of the system that combines both a
first stage detector and Viyog been done? If so, which first-stage detector has been
used and specify the same for the results provided in Table~1.''}

\rr Yes --- we now provide a complete end-to-end two-stage cascade (Energy gate
$\to$ Viyog) evaluated over 17 architectures $\times$ 10 seeds at a matched
5\% ID-FPR. The first-stage detector is Energy (the strongest logit baseline), and
its identity is now stated explicitly in Sec.~IX-A and Table~V.
\done{Table~V \& Fig.~15 (p.~11); Sec.~IX-A (p.~10, l.~569)}

% ------ DC 2 ------------------------------------------------------------------
\rc{DC~2. ``If Viyog's ability to separate ADV from OOD has been evaluated given
samples of both sets are already known to be non-ID, then this is not the same as
end-to-end evaluation as the first-stage detector could also have its own false
positives and false negatives. How does Viyog handle such situations?''}

\rr \emph{False-positive propagation is now measured directly.} When stage-1
wrongly flags an ID input, Viyog escalates it to the ADV class only
\textbf{5.2\%} of the time (median 0; exactly 0 on 14/17 architectures) versus
51--82\% for logit/\Linf{} baselines --- Viyog \emph{mitigates} rather than
amplifies upstream errors. ADV recall: adding the first-conv dormant-band
deviation to the gate lifts end-to-end ADV recall from \textbf{0.16 to 0.84}
(16/17 architectures; ViT-B the lone outlier at 0.21).
\done{Table~V \& Fig.~15 (p.~11); Sec.~IX-A (p.~10, l.~569)}

% ------ DC 3 ------------------------------------------------------------------
\rc{DC~3. ``Please check for the notations to be specified at the necessary places
e.g.: line 90, 139, and parts of the formal proposition. And please check for any
typos e.g.: Fig.~5 caption.''}

\rr We performed a global notation pass: every symbol ($\mathcal{B}$, $\nu_c$,
$\rho_c$, $r_c$, $\widetilde{\mathrm{TV}}$, T1/T2/T3) is now defined at first
use; figure captions are self-contained; the Fig.~5 caption typo is corrected.
\done{Sec.~III-A (p.~3, l.~207); symbols defined at first use; self-contained captions}

% ------ DC 4 ------------------------------------------------------------------
\rc{DC~4. ``Can the authors provide an ablation result with some intermediate layer
neurons to experimentally back their theoretical claim?''}

\rr Complete depth sweep across 6 architectures: the dormant-band \emph{shape}
statistic is \textbf{strongest at depth~0} (mean T2~0.965) and decays
monotonically, while the raw norm is \emph{weakest} there (0.52) and only peaks
mid-depth at 0.80. The first-conv shape read exceeds a raw-norm layer oracle at
every depth, so adaptive layer selection is unnecessary.
\done{Fig.~10 (p.~8); Sec.~VI-B (p.~7, l.~427)}

% ------ DC 5 ------------------------------------------------------------------
\rc{DC~5. ``Though it is acknowledged that the authors want a bimodal distribution
where `clearly anomalous inputs' get confident scores, this could lead to loss of
degree of anomaly information. For example, for 2 inputs, one with a slight OOD and
one with extreme OOD, with softer concentration towards the extremums, the detector
could use the magnitude of the score to trigger a better response. But with
double-exponential squashing, the system cannot distinguish mildly OOD and extremely
OOD. What are the comments of the authors in this regard?''}

\rr \emph{Addressed by redesign.} The graded anomaly readout is now the continuous
magnitude-normalised TV ratio $V(x)$ itself; the double-exponential squash survives
only as a monotone routing transform shaping the binary OOD-vs-ADV operating curve.
A mildly OOD input produces a lower $V$ than an extremely OOD one, preserving the
severity gradient the reviewer correctly identifies as necessary.
\done{Sec.~IV-B (p.~4, l.~336); abstract (p.~1)}

% ------ DC 6 ------------------------------------------------------------------
\rc{DC~6. ``In Fig.~2, the separation between the mean of norms of OOD and ADV is
clearly visible for the $L_\infty$ norm as stated by the authors. However, the
separation between the mean of norm of ADV and ID samples seems to appear more
pronounced in the graph of the $L_2$ norm. What do the authors comment on this?
Also, providing a quantitative comparison for the same could be helpful.''}

\rr \emph{Addressed.} We ran a 37-statistic battery at the first conv: every
raw magnitude norm sits near chance ($L_\infty$ 0.50--0.53, $L_2$ 0.58, $L_1$
similar), while the dormant-band \emph{shape} family tops $\approx$\textbf{0.98}
(Fig.~4). No magnitude norm separates ADV reliably; the deployed Viyog reads channel
\emph{shape}, not magnitude, making the $L_\infty$ vs.\ $L_2$ comparison moot ---
both are near chance for the security task T2.
\done{Fig.~4 (37-statistic battery, p.~5); Sec.~IV-B (p.~4, l.~336)}

% ------ DC 7 ------------------------------------------------------------------
\rc{DC~7. ``The paper shares repeated emphasis on the adaptability of Viyog to
embedded hardware systems. Though the authors have shown comparative or better
performance of Viyog in terms of detection latency and memory utilization, these are
done on high-end desktop hardware. No profiling on actual embedded platforms (e.g.:
Raspberry Pi, Jetson, FPGA) is provided. Can these results be extrapolated as such
to embedded targets? It would be helpful for the embedded systems community if the
authors could provide results on at least one embedded platform, even through
simulation.''}

\rr We deliver \emph{measured} costs on an H200 GPU with CUDA event timing:
Viyog runs \textbf{0.05\,ms/img}, $6\times$ faster than logit detectors
(MSP/Energy) and $347\times$ faster than MC-Dropout (30 passes); state is
\textbf{0.23\,KB} vs.\ 9--28\,MB for distance baselines (Fig.~12). A
compute-bound roofline projection for edge accelerators (Jetson Orin Nano,
RPi5+Hailo-8L) places the first-conv energy at \textbf{19--37\,$\mu$J} vs.\
0.6--1.2\,mJ for full inference (Fig.~13). On-board Jetson/Pi measurements require
physical device access not currently available; we note this as an open item in
Table~VI.
\done{Figs.~12--13 (p.~10); Sec.~VI-C \& Table~IV (p.~8, l.~451)}

% ═══════════════════════════════════════════════════════════════════════════════
\section*{Reviewer B}
\noindent\textit{We address all 12 detailed
comments (DC~1--12) below.}

% ------ DC 1 ------------------------------------------------------------------
\rc{DC~1. ``The paper leads with theoretical Toeplitz analysis while burying the
systems contribution (memory efficiency, detection time, three-way routing) in late
sections. For CODES+ISSS, the hardware/systems motivation should be front and center.
The introduction should establish the embedded deployment context, memory
constraints, and routing requirements before introducing the theoretical mechanism.''}

\rr Abstract now opens with the deployment framing (0.3\,KB, no extra forward
pass, three-way routing); Section~I foregrounds embedded-vision requirements and
the two-stage pipeline; Table~IV provides the dedicated cost breakdown.
\done{Abstract \& Sec.~I (p.~1, l.~27); Table~IV (p.~8)}

% ------ DC 2 ------------------------------------------------------------------
\rc{DC~2. ``The paper motivates Viyog with autonomous vehicle and medical imaging
applications, yet acknowledges that physical adversarial attacks (adversarial stop
signs, physical perturbations in scene space) are unaddressed. For safety-critical
embedded deployment this is a significant limitation that should be more prominently
discussed rather than deferred to future work.''}

\rr A dedicated Limitations section (Table~VI) explicitly enumerates physical/ISP
attenuation as a hard boundary condition. We note that the TV-shape mechanism
\emph{requires} digital gradient structure; physical attacks that pass through a
camera ISP low-pass filter are acknowledged as out of scope with clear reasoning.
\done{Sec.~IX-B (p.~11, l.~586); Table~VI (p.~11)}

% ------ DC 3 ------------------------------------------------------------------
\rc{DC~3. ``Results are reported from single training runs. Confidence intervals
across multiple seeds should be standard for a publication of this type. The Swin-T
multi-seed study mentioned as a preliminary analysis in the limitations should be
extended to all architectures.''}

\rr \emph{Full 20-architecture training-seed CIs delivered.} Each of the 20
CIFAR-100 backbones was re-finetuned with up to 3 independent seeds; T2
per-architecture std $\le$\textbf{0.004} across all 20 (an order of magnitude
below the cross-architecture spread); T3 std $\le$0.06. Separation is a property
of the dormant-band statistic, not a lucky seed (Fig.~8).
\done{Fig.~8 (p.~7); Sec.~VI-A (p.~6, l.~407)}

% ------ DC 4 ------------------------------------------------------------------
\rc{DC~4. ``Proposition~1 presents 2 results. (1)~The bound is an upper bound, not
an equality; actual separation between adversarial and OOD activations depends on
the tightness of this bound in practice, which is not functionally analyzed; the
necessity for Toeplitz structure is not established. (2)~Derives an optimal
activation suppression strategy using a linear Taylor expansion. This result is
attack-specific, applying to gradient sign attacks. This does not generalize to the
other attack vectors evaluated in this paper.''}

\rr \emph{Addressed by reframing the theory.}
\textit{(i)}~Toeplitz necessity removed: $V(x)$ is a per-channel shape ratio with
no structural assumption; the bound in Appendix~A holds for any affine first layer.
\textit{(ii)}~The Young envelope $\mathrm{TV}(e_c)\le\|w_c\|_1\overline{\|\nabla\delta\|_1}$
is stated as a $\le$ (empirically $\approx2\times$ loose, measured) and is not the
detection mechanism; a two-sided bound is now proven via TV subadditivity (P2).
\textit{(iii)}~The per-channel residue $e_c=w_c\!\ast\delta$ is \emph{exact by
affineness} (no Taylor, no $\mathrm{sign}\nabla\mathcal{L}$), so the
attack-agnostic property of P2 holds for FGSM, PGD, CW, and DeepFool identically.
\done{Proposition~1 (pp.~3--4); Appendix~A (p.~12, l.~623)}

% ------ DC 5 ------------------------------------------------------------------
\rc{DC~5. ``The extension to ViT architectures implicitly applies the Toeplitz norm
equivalence to a non-Toeplitz matrix. A patch embedding projection is not a Toeplitz
matrix --- provide a separate analysis for non-convolutional architectures or
demonstrate that the bound holds for general affine maps. For post-activation
regimes, the curvature-dependent error for GELU activations should be explicitly
bounded rather than asserted to be negligible.''}

\rr The theory now states the first layer as \emph{any} affine $f_0(x)=W\!\ast x+b$
(strided conv or patch-embed projection); the exact residue $e_c=w_c\!\ast\delta$
and TV-seminorm subadditivity hold identically for both without invoking Toeplitz
structure. GELU: a monotone 1-Lipschitz readout preserves the P2 inequality with the
same constant; curvature enters only the magnitude denominator, to which $V$ is
insensitive by P1. The Taylor remainder is explicitly bounded in Appendix~A
($\le0.399\,\|W\|_\infty^2\varepsilon^2$ per coordinate).
\done{Proposition~1 (pp.~3--4); Appendix~A GELU error (p.~12, l.~623)}

% ------ DC 6 ------------------------------------------------------------------
\rc{DC~6. ``There seems to be a claim of universal convergence to an interval
$(\pm 1)$ around the expected value of ID inputs in the introduction. It is walked
back later in the paper. That seems spurious and unproven.''}

\rr \emph{Addressed.} The $\pm1$ convergence claim was an assertion about the hard
double-exponential squash on the raw \Linf{} statistic, which is no longer
deployed. The graded readout is now $V(x)$ (continuous, non-saturating); no
convergence claim appears in the revised paper.
\done{Sec.~I \& abstract (p.~1, l.~27); squash demoted to routing only}

% ------ DC 7 ------------------------------------------------------------------
\rc{DC~7. ``Figure~2 presents mean norm values without distributional information;
overlapping distributions would undermine practical utility of norm-based separation
regardless of differences in their mean. The claim that `the attacker concentrates
on the single dominant channel' does not hold for FGSM and PGD, which distribute
perturbation by design.''}

\rr The original Fig.~2 (mean norm bar chart) is \textbf{replaced by the new
Fig.~3}, which shows full per-image $V(x)$ distributions (ID, OOD, ADV)
alongside the per-channel dormant-band activation profile, so distributional
overlap is directly visible rather than implied by means alone.
The ``single-dominant channel'' wording is replaced by the broadband A2 premise:
the residue $e_c=w_c\!\ast\delta$ spreads across \emph{all} channels by design,
and it is this spread that $V$ exploits.
\done{Fig.~3 (replaces orig.\ Fig.~2, p.~4); Proposition~1 premise A2 (pp.~3--4)}

% ------ DC 8 ------------------------------------------------------------------
\rc{DC~8. ``Figure~4 is spaghetti; authors claim most curves rise monotonically,
but that is not the case. Several do, but their point is undermined by the ones
that do not.''}

\rr Fig.~4 (depth sweep, now Fig.~10) shows per-arch mean $\pm$ std bands rather
than individual curves. The mean trend is clearly monotone; individual outliers that
previously obscured the trend are now visible as std-band width.
\done{Fig.~10; revised caption (p.~8)}

% ------ DC 9 ------------------------------------------------------------------
\rc{DC~9. ``Figure~5 is missing a reference to a section.''}

\rr Fig.~5 caption now cites Section~IV-B explicitly.
\done{Fig.~5 caption (p.~6)}

% ------ DC 10 -----------------------------------------------------------------
\rc{DC~10. ``In Figure~7 authors drop MCD from the detection time figure on
visualization grounds rather than using a log-scale axis or separate table entry;
this is a questionable presentation choice that obscures a meaningful data point.''}

\rr MCD ($347\times$ slower than Viyog) is retained in Fig.~12(a) on a log-scale
axis. The caption notes that it was omitted from the linear-scale scatter to
preserve resolution for the other methods, not to hide the data.
\done{Fig.~12(a) log-scale; caption note (p.~10)}

% ------ DC 11 -----------------------------------------------------------------
\rc{DC~11. ``62 references for a 9-page paper leaves insufficient space for thorough
method exposition and suggests the related work section may be disproportionate
relative to the technical contribution.''}

\rr The revised paper fits within the 14-page limit (expanded scope and
experiments); the reference count is now proportionate to the expanded coverage.
Redundant citations have additionally been pruned from 62 to 47.
\done{Reference list (pruned 62$\to$47; \texttt{final.tex} bibliography, pp.~13--14)}

% ------ DC 12 -----------------------------------------------------------------
\rc{DC~12. ``The writing style exhibits inconsistencies including irregular em-dash
usage and unexplained italicization that suggest uneven editorial attention. The
authors should revise for stylistic consistency.''}

\rr A global style pass unified em-dash usage, removed unexplained italicisation,
and standardised all notation throughout.
\done{Global notation/style pass (throughout final.tex; em-dashes, italics, symbol definitions at first use; throughout)}

% ═══════════════════════════════════════════════════════════════════════════════
\section*{Reviewer C}
\noindent\textit{Reviewer C's Detailed Comments
section restates the same content as the weaknesses; we address all four weakness
points (W1--W4) below.}

% ------ W1 -------------------------------------------------------------------
\rc{W1. ``The paper shows that a white box attacker can bypass Viyog by adding a
simple norm preservation penalty to the attack. This suggests that the method can be
defeated with small changes to standard PGD. The authors describe this failure mode
as safer, but it still raises concerns about robustness.''}

\rr \emph{Strengthened to 0\% evasion.} Proposition~1 (P1) proves $V$ is
$0$-homogeneous per channel: a norm-preserving attacker \emph{cannot move} $V$ by
construction. Empirically, on 20 architectures (Fig.~14):
\begin{itemize}[leftmargin=1.2em,nosep]
\item A0 (standard PGD): $V$ AUROC~0.93.
\item A1 (norm-preserving): \textbf{0\% evasion} --- P1 holds exactly.
\item A2 (TV-aware, white-box on deployed $V$): erodes to 0.64 worst-case;
      the complementary HF read holds at 0.74.
\item A3 (both-aware): costs the attacker \textbf{8\%} success (A2~88.9\%$\to$A3~81.1\%);
      the combined $\max(V_{\mathrm{TV}},V_{\mathrm{HF}})$ floors at AUROC~0.68.
\item EOT stochastic defence: restores $V$ to \textbf{0.82}; 18/20 archs
      improve above the diagonal.
\end{itemize}
On GTSRB the combined detector floors at \textbf{0.775} while attack success
$\ge$80\% --- an attacker either keeps the detector above chance or surrenders
most success on both datasets.
\done{Proposition~1 (pp.~3--4); Sec.~VII A0--A3$+$EOT (p.~8); Fig.~14 (p.~11)}

% ------ W2 -------------------------------------------------------------------
\rc{W2. ``The CNN results are reported as single seed values. The authors
acknowledge this limitation, but multi seed evaluation is only provided for Swin-T.
This leaves some uncertainty about the stability and reproducibility of the main
results.''}

\rr Addressed above (B~DC~3): full 20-arch training-seed CIs are now delivered.
Per-architecture T2 std $\le$0.004 --- the headline separation is a property of
the dormant-band statistic, not a lucky seed.
\done{Fig.~8 (p.~7)}

% ------ W3 -------------------------------------------------------------------
\rc{W3. ``The lowest AUROC reported is 75.53 on DenseNet with MNIST, which is a
relatively simple dataset. Also, all experiments are limited to image classification.
It is not clear how well the method would work for NLP or tabular data.''}

\rr \emph{Two separate issues resolved.}
(i)~The DenseNet-MNIST degeneracy (25/64 dead first-conv filters) is fixed by
active-channel selection, restoring T2 to 0.80.
(ii)~75.53 is a \emph{different} evaluation (DenseNet, MNIST-as-OOD, task T1);
we list it as a genuine limitation.
(iii)~Non-vision is now probed with a 1D-CNN on band-limited multi-sinusoid
signals: Viyog T2 \textbf{0.96} vs.\ raw \Linf{} 0.71 at the first
\texttt{Conv1d} (3 seeds, Fig.~18), confirming the mechanism is modality-agnostic.
\done{Sec.~IX-B limitation (ii) (p.~11, l.~586); App.~B \& Fig.~18 (p.~13)}

% ------ W4 -------------------------------------------------------------------
\rc{W4. ``Overall, the proposed technique is vulnerable to adaptive attacks and
lacks multi seed evaluation for the main results, which limits confidence.''}

\rr Both concerns are now resolved: 0\% evasion under norm-preserving A1 (P1),
full A0--A3 ladder showing the cost-of-evasion trade-off (Sec.~VII), and
full-20-arch seed CIs (Fig.~8). See W1 and W2 above for details.
\done{Fig.~8; Fig.~14; Proposition~1 (pp.~3--4, 7, 11)}

% ═══════════════════════════════════════════════════════════════════════════════
\section*{Reviewer D}
\noindent\textit{We address all 3 detailed comment
paragraphs (DC~1--3) below.}

% ------ DC 1 ------------------------------------------------------------------
\rc{DC~1. ``An important improvement would be to evaluate Viyog as part of a
complete anomaly-handling pipeline rather than in isolation. Integrating it with
common first-stage OOD detectors and reporting system-level metrics would clarify
how errors propagate across stages. This would help quantify real-world reliability
and identify whether Viyog mitigates or amplifies mistakes made upstream.''}

\rr This concern is fully addressed by the same two-stage cascade experiment
described above in A~DC~1--2. Viyog \emph{mitigates} upstream
FP errors: end-to-end ADV recall lifts from 0.16 to 0.84; ID$\to$ADV
mis-escalation drops to 5.2\% (median 0).
\done{Table~V \& Fig.~15 (p.~11); Sec.~IX-A (p.~10, l.~569)}

% ------ DC 2 ------------------------------------------------------------------
\rc{DC~2. ``The method could be improved by explicitly addressing near-distribution
or semantically similar OOD inputs. Evaluations on harder OOD benchmarks, combined
with augmenting the detector using complementary early-layer statistics or adaptive
layer selection, would strengthen robustness in ambiguous scenarios where OOD inputs
resemble in-distribution data.''}

\rr \emph{All three requests addressed.}
\emph{Near-OOD split (17 archs):} far-OOD T3~\textbf{0.96}, near-OOD~\textbf{0.74}
(up to 0.95), texture-OOD~0.67 --- graceful degradation, not collapse; a new
figure (Fig.~11) plots the full far/near/texture breakdown and shows a
high-frequency dorm-band read \emph{complements} the total-variation read,
recovering texture-OOD to \textbf{0.92}.
\emph{Complementary statistics:} logit+Viyog gives balanced accuracy
\textbf{0.757/0.768} (CIFAR-100/GTSRB); the full LDA panel reaches
\textbf{0.798/0.863}; bootstrap 95\% CIs are disjoint from Energy-only (Fig.~7).
\emph{Adaptive layer selection:} depth sweep confirms first conv is the global
optimum for $V$ (mean T2~0.965), exceeding a raw-norm layer oracle at every depth.
\emph{T3 breakdown across 20 archs, 4 attacks, 3 datasets} (Fig.~9): minimum T3
across all conditions is 0.77; FGSM shows the highest T3 (0.89) and APGD-CE the
lowest (0.77), consistent across all three datasets.
\done{Figs.~7, 9, 10, 11 (pp.~7--9); Secs.~VI-A (p.~6, l.~407), VI-B (p.~7, l.~427)}

% ------ DC 3 ------------------------------------------------------------------
\rc{DC~3. ``Viyog would benefit from a more comprehensive adaptive threat model.
Testing against multi-objective attacks that simultaneously preserve norms and
manipulate early activations, along with simple defense-hardening strategies such as
stochastic feature sampling, would better characterise worst-case robustness and
reinforce the safety claims.''}

\rr \emph{Completed.} EOT evaluation across 20 architectures (Fig.~14b):
the stochastic dorm-band defence restores mean T2 to \textbf{0.82} where a
fixed-band read drops to 0.68 under the same attacker; 18/20 architectures
improve above the diagonal.
A stronger held-out-basis A4 attacker (6 archs) reaches the same qualitative
conclusion with $\approx$15\% success surrendered.
\done{Fig.~14(b) EOT panel (p.~11); Sec.~VII-B (p.~9, l.~481)}

% ── Closing summary ──────────────────────────────────────────────────────────
\noindent\rule{\columnwidth}{0.4pt}

\noindent\textbf{Completeness of new experiments.} The revised manuscript now
reports experiments across \textbf{3 datasets $\times$ 20 architectures $\times$
4 standard attacks $\times$ 11 detectors} plus the adaptive families:
full-20-architecture training-seed CIs, EOT (20 archs), measured H200 latency
(all 11 \texttt{pytorch-ood} baselines), an on-silicon roofline (two edge
accelerators), a 1D non-vision sanity check, the near/texture-OOD breakdown, and
the A0--A3 adaptive ladder on two datasets.

\noindent\rule{\columnwidth}{0.4pt}

\vspace{3pt}
\noindent We believe the revised manuscript fully addresses all concerns and are
grateful for the detailed, constructive reviews.

\balance
\end{document}
